\section{Prosjektforsvar: \emph{Coop App-prosjektet i lys av pensum}}
\label{sec:project_defense}
\settermlabel{terms:lec09}

\subsection*{Oversikt}
Dette kapittelet kobler eksamensprosjektets innhold (rådgivning til Thor
Skov Jørgensen om Coop App, levert 15.05.2026) til pensum. Formålet er å
forberede forsvar -- både ved å vise \emph{hvorfor} vi valgte de rammene vi
brukte (DVC og DBM), og ved å peke på hvordan resten av pensum styrker eller
nyanserer argumentet. Det kapittelet \emph{ikke} gjør, er å bytte ut
prosjektets konklusjoner; det \emph{utvider} det teoretiske ankret bak dem.

\subsection{Prosjektets kjerne i én setning}
\keypoint{Coop App og Lobyco er strategisk verdifulle bare i den grad de styres
mot dokumentert butikk- og medlemsverdi -- ikke aktivitetstall. Anbefalingen
er videreføring med tre styringsdimensjoner: KPI-hierarki, kjedespesifikk
diversifisering innenfor felles plattform, og Lobyco/retail media som
Coop-først-kapabilitet.}

\subsection{De tre anbefalingene og deres teorianker}

\begin{tabular}{p{3.5cm}p{6cm}p{4.5cm}}
\hline
\textbf{Anbefaling} & \textbf{Argumentet} & \textbf{Primæranker} \\
\hline
1. KPI-hierarki &
Skille bruksmål, atferdsmål og lønnsomhetsmål; styre etter de to siste. &
DVC outputs $\neq$ outcomes (Duus \& Cooray 2020); DBM value
\emph{capture} (Weill \& Woerner 2018). \\
\hline
2. Kjedespesifikk\newline diversifisering &
Moduler innenfor én felles plattform; 365discount, Kvickly, SuperBrugsen
har ulike kundebehov. &
DVC \emph{relevance} som kjerne; DBM omnichannel-konsistens. \\
\hline
3. Lobyco \& retail media som Coop-først &
Klar governance; retail media testes med eksplisitte krav. &
DBM modular-producer-/omnichannel-spenning; Bjørn-Andersen
\& Hedman (2016) om legacy debt. \\
\hline
\end{tabular}

\subsection{Hvorfor DVC + DBM? (forsvar av rammeverk-valget)}

\subsubsection*{Rammeverkene som komplementært par}
Rammeverkene er valgt fordi de svarer på \emph{to ulike spørsmål}:
\begin{itemize}
  \item \term{DVC} (Duus \& Cooray 2020) svarer på: \emph{hvor og hvordan kan
    digital kundeverdi skapes?} Den bryter ned digital verdi i experiences,
    relationships, evolution, med relevance som kjerne -- omgitt av et ytre
    lag av digital infrastruktur, kompetanse og outputs.
  \item \term{DBM} (Weill \& Woerner 2018) svarer på: \emph{passer den
    digitale verdien inn i forretningsmodellen, og kan den \emph{fanges}?}
    De fire arketypene (Supplier, Omnichannel, Modular producer, Ecosystem
    driver) plasserer virksomheten i kvadranter etter \emph{customer
    knowledge} og \emph{business design}.
\end{itemize}

\keypoint{Den analytiske brutaliteten i å bruke dem sammen: DVC alene kan
overvurdere digital aktivitet; DBM alene kan undervurdere det digitale
potensialet. Sammen tvinger de fram skillet mellom value creation og value
capture -- selve prosjektets røde tråd.}

\subsubsection*{Styrker og svakheter (klar til å forsvare)}
\begin{itemize}
  \item \textbf{DVC styrke:} gir konkret språk for hvor digital verdi
    oppstår.
  \item \textbf{DVC svakhet:} kan gjøre aktivitet til å \emph{virke}
    verdifull bare fordi den eksisterer -- mangler i seg selv finansiell
    disiplin (jf.\ \autoref{sec:lec03}).
  \item \textbf{DBM styrke:} tvinger value-capture-spørsmålet; gir
    arketyper som strategisk språk.
  \item \textbf{DBM svakhet:} virkelige virksomheter passer ikke pent i én
    boks. Coop er overveiende \emph{omnichannel}, men Lobyco trekker mot
    \emph{modular producer}. Rammeverket er beskrivende, ikke
    foreskrivende.
\end{itemize}

\subsection{Hvor resten av pensum styrker forsvaret}

Selv om prosjektet bruker DVC + DBM som hovedrammer, kan nesten alle
forelesningene mobiliseres til å \emph{forsvare} de tre anbefalingene. Her
er pensum-koblingene en eksaminator vil sette pris på:

\subsubsection*{\autoref{sec:lec01} -- Porter (1996): Strategi som
posisjonering}
Anbefalingene støttes av Porters skille mellom \term{strategi} og
\term{operasjonell effektivitet}. Coops butikk-først omnichannel er en
\emph{posisjon}; appen må forsterke den, ikke konkurrere med den.
Anbefaling~1 (KPI-hierarki) er i realiteten Porters \emph{trade-off}-logikk
oversatt til måling: alt kan ikke optimaliseres samtidig, og vi velger å
optimalisere butikkøkonomi -- bevisst nedprioritering av leverandørtall som
hovedmål.

\subsubsection*{\autoref{sec:lec02} -- Five Forces, verdikjede, VRIO, SWOT}
\begin{itemize}
  \item \textbf{Five Forces:} Coop opererer i en moden bransje med høy
    rivalisering, sterk leverandørmakt, og økende digital substitusjonstrussel
    (Nemlig, Lidl, plattformaktører).
  \item \textbf{VRIO:} Coops 1{,}8M medlemsdata er en kandidat til
    \emph{bærekraftig konkurransefortrinn} -- verdifullt, relativt sjeldent
    (kun Salling og Coop har sammenlignbare lojalitetsprogrammer i Danmark),
    vanskelig å imitere uten samme medlemsbase. \emph{O}-leddet
    (organisert utnyttelse) er nettopp det Anbefaling 1 og 3 sikrer.
  \item \textbf{Value chain:} Coop App og Lobyco er teknologi-/markedsførings-
    støtteaktiviteter som må styrkes primæraktiviteten (butikk).
\end{itemize}

\subsubsection*{\autoref{sec:lec04} -- Wessel et al.\ (2021): DT vs.\ ITEOT}
\keypoint{Dette er det kraftigste teoritillegget for forsvar.} Prosjektet
argumenterer implisitt for at Coop driver \emph{IT-Enabled Organizational
Transformation}, ikke \emph{Digital Transformation}. Bevisene:
\begin{itemize}
  \item Coop redefinerer ikke sin identitet -- de er fortsatt en
    butikkbasert dagligvare-konsern.
  \item Lukkingen av Coop.dk (januar 2025) er det sterkeste empiriske
    beviset for ITEOT-tolkningen -- en DT-virksomhet ville beveget seg
    \emph{mot} digital salg, ikke trukket seg ut av det.
  \item Value proposition forsterkes (lettere butikkbesøk via Scan \& Pay),
    ikke redefineres.
\end{itemize}
\textbf{Konsekvens for anbefalingene:} en ITEOT-tolkning gjør Anbefaling~3
(Coop-først-styring av Lobyco) teoretisk obligatorisk. Hvis Coop skulle
drive DT mot \emph{ecosystem driver}-identitet, ville Lobyco vært
strategisk kjerne. Siden Coop driver ITEOT, må Lobyco styres som støtte for
den eksisterende identiteten.

\subsubsection*{\autoref{sec:lec05} -- Digital innovasjon og disrupsjon}
\begin{itemize}
  \item Coop App er en \term{sustaining innovation} (forbedrer eksisterende
    verdiforslag for eksisterende kunder), ikke en disruptiv innovasjon.
  \item Samtidig er den \term{digital innovasjon (DgI)} i Yoo et al.\ sin
    forstand: en \emph{rekombinasjon} av digitale (kundedata, push,
    Scan-betaling) og fysiske (butikkbesøk, sortiment) komponenter.
  \item De \emph{potensielle} disruptorene er presise: rene
    netthandelsaktører (low-end på service), lavpriskjeder (low-end på
    pris), og plattformaktører (new-market). Anbefalingen om
    diversifisering (2) møter low-end-pris-disrupsjon ved å tilpasse
    365discount-modulen til pris-/tilbudsfokus.
\end{itemize}

\subsubsection*{\autoref{sec:lec06} -- Økosystemer og plattformer}
\begin{itemize}
  \item Coop er \emph{ikke} en ecosystem driver -- og prosjektet er
    eksplisitt på at det er korrekt diagnose. Ingen evidens for
    plattformidentitet eller bred ekosystemkoordinering.
  \item Lobyco er en potensiell \term{modular producer} i \emph{andres}
    plattformer/økosystemer. Anbefaling~3 handler nettopp om å håndtere
    spenningen mellom Coops omnichannel-rolle og Lobycos
    modular-producer-rolle.
  \item \term{Network effects} på Coop App er overveiende
    \emph{single-sided} (medlemmer på samme side). Retail media legger til
    en \emph{cross-side}-effekt mellom medlemmer og leverandører -- men den
    cross-side-effekten kan bli \emph{negativ} for medlemmene hvis
    leverandørstyrt innhold fortrenger relevans. Anbefaling 3 adresserer
    dette eksplisitt.
\end{itemize}

\subsubsection*{\autoref{sec:lec07} -- Implementering og legacy debt}
\begin{itemize}
  \item \textbf{Bjørn-Andersen \& Hedman (2016)} er sentralt evidensanker for
    Anbefaling 1 (KPI-hierarki) og 3 (Lobyco-styring). Coop har historisk
    erfart at IT-investeringer uten klar effekt og uten støtte fra ledelsen
    har gitt teknologisk gjeld.
  \item \textbf{Tawse \& Tabesh (2021)} gir et språk for hva som faktisk skal
    til for å gjennomføre anbefalingene:
  \begin{itemize}
    \item \emph{Competency} -- analytisk kapasitet til å bygge og bruke
      KPI-hierarkiet.
    \item \emph{Commitment} -- ledelsens støtte (Skov Jørgensens
      mandatering).
    \item \emph{Coordination} -- mellom Coop-konsernet, kjedene og Lobyco.
  \end{itemize}
  \item \textbf{Miller (1997):} prosjektets implisitte risikoavveiing svarer
    til Millers \emph{propitiousness} -- man styrer etter dokumentert verdi
    fordi omstendighetene (lav margin, økonomisk press) krever det.
  \item \textbf{Hybrid governance} er den implisitte styringsmodellen: noen
    beslutninger sentralt (plattform, KPI-rammeverk), noen lokalt (kjedenes
    moduler).
\end{itemize}

\subsection{Sannsynlige forsvarsspørsmål -- og svar}

\subsubsection*{\guillemotleft{}Hvorfor disse to rammeverkene -- DVC og DBM?\guillemotright{}}
De svarer på to ulike spørsmål: \emph{hvor} digital verdi kan skapes (DVC) og
\emph{om} den fanges (DBM). Sammen forhindrer de to typiske feil --
overvurdering av aktivitet (DVC alene) og undervurdering av digitalt
potensial (rent finansiell linse). De er begrepsmessig komplementære, ikke
overlappende.

\subsubsection*{\guillemotleft{}Hvorfor ikke Porter / Wessel / Christensen som hovedlinse?\guillemotright{}}
Porter er underliggende premiss for hele oppgaven (strategisk
posisjonering), men gir ikke konkret digital diagnostikk. Wessels DT/ITEOT
er kraftig, men anvendes implisitt -- Coop er ITEOT. Christensen er relevant
for å identifisere \emph{trusler} (disruptorer), men ikke for å vurdere
appens egen verdi. DVC + DBM er valgt fordi de \emph{operasjonaliserer} det
digitale verdispørsmålet i en form som direkte gir anbefalinger.

\subsubsection*{\guillemotleft{}Hvor i Weill \& Woerner-matrisen er Coop?\guillemotright{}}
\textbf{Omnichannel}, entydig. Bevis: høy customer knowledge via medlemsdata,
business design er value-chain-orientert (Coop eier butikkene, koordinerer
sortiment), Coop.dk-lukkingen utelukker ecosystem-driver-ambisjonen.
\emph{Lobyco} trekker isolert mot modular producer, men er kapabilitet i
Coop-systemet -- ikke selvstendig strategisk posisjon.

\subsubsection*{\guillemotleft{}Hvorfor ikke selge Lobyco?\guillemotright{}}
Kildegrunnlaget er ikke tilstrekkelig. Vi kjenner ikke Lobycos kostnader,
eksterne inntekter, bokverdi eller hvor avhengig Coop App er av
Lobyco-teknologien. Å anbefale salg basert på et kommersielt-tilfarget
offentlig datasett ville være intellektuelt uærlig. Anbefaling~3 (klar
styring + Coop-først-mandat) er den intellektuelt forsvarlige posisjonen
gitt usikkerhetsnivået.

\subsubsection*{\guillemotleft{}Hva er det største kildekritiske problemet?\guillemotright{}}
\keypoint{Evidenskjeden er asymmetrisk dokumentert.}
\begin{enumerate}
  \item Adoption/rekkevidde (Lobyco: 1{,}8M brukere) -- godt dokumentert,
    men kommersiell kilde.
  \item Engasjement/atferdsmål (Playable: gamification-effekter) -- delvis
    dokumentert.
  \item Butikkøkonomi (kurvstørrelse, trafikk, margin tilbakeført til
    appen) -- \emph{ikke} dokumentert i offentlig kildegrunnlag.
  \item Lønnsomhet / strategisk fortrinn -- \emph{ikke} dokumentert.
\end{enumerate}
Anbefaling 1 (KPI-hierarki) er nettopp svaret på dette: konsernet trenger
\emph{eget} datagrunnlag for trinn 3 og 4.

\subsubsection*{\guillemotleft{}Hva er seleksjonsbias-problemet?\guillemotright{}}
Appbrukere kan være de mest lojale kundene \emph{før} de laster ned appen.
Hvis de handler oftere, kan korrelasjonen reflektere prior lojalitet, ikke
appens effekt. Lobyco-tallet \guillemotleft{}appkunder handler 50~\% oftere\guillemotright{} er bias-
sårbart. Uten randomisert tilordning eller naturlige eksperimenter (f.eks.\
ulik utrullingstidsplan på tvers av regioner) kan kausalitet ikke
påstås. Dette er en \emph{begrensning ved analysen}, ikke et argument mot
appen.

\subsubsection*{\guillemotleft{}Hvorfor kjedespesifikk diversifisering og ikke én felles
app?\guillemotright{}}
DVC-relevance er ikke uniform på tvers av kundebasen. 365discount-kunder
prioriterer tilbud og enkelhet; Kvickly/SuperBrugsen-kunder prioriterer
inspirasjon og bredt sortiment. DBM-omnichannel-modellen krever at den
digitale relasjonen \emph{reflekterer} kjedens distinkte kundeløfte; ellers
svekkes konsistensen som er omnichannel-modellens kjerne. Moduler i felles
plattform balanserer relevance (per kjede) mot kompleksitet (felles
infrastruktur via Lobyco).

\subsubsection*{\guillemotleft{}Hvilke risikoer ved selve anbefalingene?\guillemotright{}}
\begin{itemize}
  \item \textbf{Anbefaling 1:} Bygging av måleinfrastruktur har egne
    kostnader; strengere måle-krav kan oppleves som brems på utvikling.
    Akseptabel risiko fordi alternativet er å styre blindt.
  \item \textbf{Anbefaling 2:} Modulær diversifisering øker
    innholdsproduksjon og koordinasjonskompleksitet. Risiko mitigeres ved
    felles plattform.
  \item \textbf{Anbefaling 3:} Strenge styringsrammer kan svekke Lobycos
    innovasjonsevne og redusere ekstern omsetning. Akseptabel risiko gitt
    butikk-først-mandatet.
\end{itemize}

\subsubsection*{\guillemotleft{}Hva er distinksjonen mellom value creation og value
capture i denne casen?\guillemotright{}}
\begin{itemize}
  \item \emph{Value creation:} Kunden får relevans, bekvemmelighet, bonus
    -- godt dokumentert gjennom DVC-laget (experiences, relationships,
    relevance).
  \item \emph{Value capture:} Coop tjener på det -- gjennom inkrementell
    butikktrafikk, høyere kurvstørrelse, leverandørinntekter (retail
    media). \emph{Ikke} dokumentert med Coops egne tall.
\end{itemize}
Hele prosjektets \guillemotleft{}strategic stance\guillemotright{} er at value creation $\neq$ value
capture, og at konsernets oppgave er å lukke gapet med målbar styring.

\subsubsection*{\guillemotleft{}Driver Coop digital transformation eller IT-enabled
transformation?\guillemotright{}}
\textbf{IT-enabled organizational transformation}, basert på Wessel et al.\
(2021). Coops identitet forblir \guillemotleft{}vi er et butikkbasert dagligvare-konsern\guillemotright{}.
Den digitale teknologien forsterker eksisterende value proposition, ikke
redefinerer den. Lukkingen av Coop.dk i januar 2025 er det avgjørende
empiriske beviset -- en virksomhet i DT ville beveget seg motsatt vei.
\emph{Implikasjon for anbefalingene:} ITEOT-rammen gjør Anbefaling 3
(Coop-først-styring av Lobyco) teoretisk korrekt.

\subsubsection*{\guillemotleft{}Hva ville et alternativt teorivalg gitt oss?\guillemotright{}}
\begin{itemize}
  \item \textbf{Christensen alene} -- ville fokusert på trusler og
    selv-disrupsjon. Manglet språk for å vurdere appens \emph{positive}
    bidrag.
  \item \textbf{Wessel et al.\ alene} -- ville gitt presis diagnose (ITEOT)
    men ikke verktøy for hvor digital verdi konkret oppstår.
  \item \textbf{Porters Five Forces alene} -- bransjeanalyse, ikke
    forretningsmodellanalyse.
  \item \textbf{Plattform-/økosystemlitteratur alene} -- ville pålagt Coop en
    ecosystem-driver-ambisjon kildegrunnlaget ikke støtter.
\end{itemize}
DVC + DBM er valgt fordi de operasjonaliserer det \emph{rette} spørsmålet:
hvordan styre digital aktivitet mot dokumentert butikkverdi.

\subsection{Strukturen for å svare på pensumspørsmål (CRE+I)}

Etter Lektion 9 sine eksamenstips, bruk denne strukturen:
\begin{enumerate}
  \item \textbf{Claim} -- påstand (\textit{hva du hevder}).
  \item \textbf{Reason} -- begrunnelse (\textit{hvorfor det er sant}).
  \item \textbf{Evidence} -- evidens (\textit{kilder/data som støtter}).
  \item \textbf{Innvending} -- motforestilling og ditt motsvar.
\end{enumerate}

\subsection{Begrensninger ved oppgaven (vær åpen om disse)}
\begin{itemize}
  \item Ingen interne Coop-data (margin, kostnader, inkrementalitet).
  \item Alle app/Lobyco-kilder er kommersielt påvirket (Lobyco, Playable,
    Shortcut).
  \item Seleksjonsbias i adferdsdata om appbrukere.
  \item Internasjonale sammenligningskilder (ICA, Tesco) viser korrelasjon,
    ikke kausalitet.
  \item Bjørn-Andersen \& Hedman (2016) er fra 2016 og kan ikke ukritisk
    brukes om Coops nåværende situasjon.
\end{itemize}

\subsection{En tre-setningers \guillemotleft{}elevator pitch\guillemotright{} for forsvaret}

\keypoint{Vi rådga Thor Skov Jørgensen om Coop App fordi den sentrale
strategiske utfordringen er styring, ikke investeringsnivå. Vi brukte DVC
til å vise hvor digital kundeverdi skapes, og DBM til å vise om den fanges
i Coops omnichannel-forretningsmodell -- og fant et systematisk gap mellom
skapt og fanget verdi. Anbefalingen er videreføring av appen, men med
KPI-hierarki, kjedespesifikk diversifisering og Coop-først-styring av
Lobyco, fordi disse tre styringsdimensjonene er det som tvinger digital
aktivitet til å bli butikkøkonomisk verdi.}
